\section{Formal Model}

Let $\mathcal S$ and $\mathcal A$ be finite. A length-$H$ controlled process $M$ consists of an initial law $\rho(s_1)$, transition kernels $P_t(s_{t+1}\mid z_t)$, and reward kernels $Q_t(r_t\mid z_t)$ supported on $[0,1]$. We allow full history dependence; Markov decision processes are a special case. For
$z_t=(s_1,a_1,\ldots,s_t,a_t)$ and $h_t=(s_1,a_1,\ldots,a_{t-1},s_t)$, define
\[
L_t^\eta(z_t)=\prod_{j=1}^t\eta_j(a_j\mid h_j)
\]
and the environment prefix factor
\[
g_{M,t}(z_t)=\rho(s_1)\prod_{j=1}^{t-1}P_j(s_{j+1}\mid z_j).
\]
Thus the state-action prefix mass is
$p_{M,t}^\eta(z_t)=g_{M,t}(z_t)L_t^\eta(z_t)$.
Let $m_{M,t}(z_t)=\bbE_M[R_t\mid z_t]$.

The logged distribution contains complete episodes under $\beta$. Two environments are observationally equivalent when they induce the same distribution of states, actions, and rewards under $\beta$. A contrast is uniformly identified when it has the same value in every pair of observationally equivalent environments in the model class.

\section{Proof of the Uniform Identification Theorem}

\begin{theorem}[Restated]
For fixed finite alphabets and policies $(\beta,\pi,\pi_0)$, the contrast $V_M^\pi-V_M^{\pi_0}$ is uniformly identified over all finite-horizon controlled processes with rewards in $[0,1]$ if and only if, for every formal prefix $z_t$,
\[
L_t^\beta(z_t)=0\Longrightarrow L_t^\pi(z_t)=L_t^{\pi_0}(z_t).
\]
When the condition holds,
\[
V_M^\pi-V_M^{\pi_0}
=\bbE_M^\beta\left[\sum_{t=1}^H\gamma^{t-1}
\frac{L_t^\pi(Z_t)-L_t^{\pi_0}(Z_t)}{L_t^\beta(Z_t)}R_t\right],
\]
where the ratio is evaluated only on logged prefixes.
\end{theorem}

\begin{proof}[Sufficiency]
For each $t$,
\begin{align*}
&\bbE_M^\pi[R_t]-\bbE_M^{\pi_0}[R_t]\\
&\quad=\sum_{z_t}g_{M,t}(z_t)
\{L_t^\pi(z_t)-L_t^{\pi_0}(z_t)\}m_{M,t}(z_t).
\end{align*}
Partition the sum according to whether $L_t^\beta(z_t)$ is positive. On the zero-likelihood set, the braced difference is zero by \UCC. On the positive-likelihood set, multiply and divide by $L_t^\beta(z_t)$:
\begin{align*}
&\sum_{z_t:L_t^\beta>0}g_{M,t}(z_t)L_t^\beta(z_t)
\frac{L_t^\pi(z_t)-L_t^{\pi_0}(z_t)}{L_t^\beta(z_t)}m_{M,t}(z_t)\\
&=\bbE_M^\beta[W_t^\Delta(Z_t)R_t].
\end{align*}
Summing with discount factors proves the representation. Its right-hand side is a functional of the logged law and known policies, so the contrast is identified.
\end{proof}

\begin{proof}[Necessity]
Suppose \UCC fails. Then for some $t$ and formal prefix
$z_t^\star=(s_1^\star,a_1^\star,\ldots,s_t^\star,a_t^\star)$,
$L_t^\beta(z_t^\star)=0$ but
$d^\star:=L_t^\pi(z_t^\star)-L_t^{\pi_0}(z_t^\star)\ne0$.
Construct a deterministic history-tree process whose state at depth $j$ records the preceding action sequence and whose deterministic transitions make $z_t^\star$ feasible when its action sequence is selected. All rewards are zero except possibly at $z_t^\star$. Let $M_0$ set its reward to zero and $M_1$ set it to one.

Because the logger assigns zero probability to the action prefix, it never observes $z_t^\star$. Hence $M_0$ and $M_1$ induce exactly the same logged trajectory distribution. Their contrasts differ by
$\gamma^{t-1}d^\star\ne0$. Therefore the contrast is not a functional of the logged law and is not uniformly identified.
\end{proof}

